\chapter{METODOLOGIA}
\label{capitulo-metodologia}

% Nesta parte dedique-se à apresentação dos materiais usados e da metodologia assumida.

\section{Materiais}
\label{capitulo-metodologia-materiais}
% TODO: preencher
% Descrever os materiais utilizados (de acordo com a especificação técnica), e quantidade conforme atividade a ser realizada. Mostrar ferramentas usadas, objetos, aplicativos, equipamentos, etc.
% Evite simplesmente listar os materiais, procure descrever as etapas (PROCEDIMENTOS) da atividade prática. Devem ser apresentadas as equações e fotos da atividade prática desenvolvida. Não deve conter os resultados do trabalho.  
% As figuras no trabalho devem ser anunciadas no texto, centralizada e o nome da figura deve ser colocado em cima da mesma, com letra Arial 12, espaçamento entre linhas simples. A fonte deve ser colocada embaixo da figura com letra Arial 10, espaçamento simples entre linhas. 
% Quando se colocam quadros ou tabelas, estes também devem ser centralizados. As tabelas e quadros devem ser de boa qualidade, inseridas no texto. A indicação da fonte é feita na parte inferior da tabela/quadro. As tabelas contêm números e os quadros somente texto e ambos precisam trazer descrição dos resultados, ou seja, a tabela ou quadro precisa ser explicado para o leitor.
% Todo material citado deve constar nas referências. Deve ser apresentada preferencialmente em ordem cronológica.

\section{Métodos}
\label{capitulo-metodologia-metodos}
% TODO: preencher
% O conteúdo desta parte do trabalho deve conter: a descrição da área e do universo de estudo, as fontes que serão utilizadas para a coleta de dados para realização do estudo, população amostra ou sujeitos da extensão e/ou empresa alvo do estudo, tipo e método da pesquisa usado na construção da extensão, bem como os procedimentos e métodos escolhidos para a extensão (como pretende fazer).
% Apresentar os procedimentos metodológicos do estudo ocorridos no projeto de extensão, tais como: os referentes à natureza da extensão (básica ou aplicada), à abordagem do problema (quantitativa ou qualitativa), aos objetivos (exploratória, descritiva, explicativa) e aos procedimentos metodológicos a serem adotados, todos devem ser citados e explicados, justificando que realmente se trata desse tipo de extensão que você acredita que seja.
% Apresentar também as metodologias próprias da área de estudo citando normas, leis, métodos cientificamente aceitos, como por exemplo, a implementação, estruturação, criação ou utilização de um software.

\section{Atividades e Cronograma}
\label{sec:atividades-cronograma}

A execução do projeto será dividida em etapas, contemplando levantamento de requisitos, pesquisa das tecnologias, montagem do protótipo, desenvolvimento do software embarcado, testes e apresentação dos resultados. O cronograma abaixo é uma proposta inicial e poderá ser ajustado conforme o andamento do projeto.

\begin{table}[htb]
\centering
\small
\caption{Cronograma de execução das atividades da proposta de extensão}
\label{tab:cronograma}
\begin{tabular}{|p{5.5cm}|c|c|c|c|}
\hline
\textbf{Atividade} & \textbf{Semana 1} & \textbf{Semana 2} & \textbf{Semana 3} & \textbf{Semana 4} \\
\hline
1. Pesquisa e levantamento de requisitos & X & --- & --- & --- \\
\hline
2. Projeto e montagem do bracelete & --- & X & --- & --- \\
\hline
3. Programação do ESP32 e integração dos sensores & --- & --- & X & --- \\
\hline
4. Testes, avaliação e apresentação do protótipo & --- & --- & --- & X \\
\hline
5. Documentação e registro dos resultados & --- & --- & --- & X \\
\hline
\end{tabular}
\fonte{Elaborado pelos autores (2026).}
\end{table}
